\section{Problem Formulation}
\label{sec:problem-setting}



\textbf{From AI for AI to meta-evolution.}
AI for AI describes the object of optimization: AI systems participate in creating or improving other AI systems, including training code, model architectures, algorithms, agent harnesses, and AI hardware.
Evolution describes the optimization process: an AI repeatedly modifies a candidate system according to feedback from its execution.
Meta-evolution closes an additional learning loop by reusing these evolutionary trajectories to train the model that proposes future modifications.
Recursive self-improvement (RSI) requires a stronger and sustained loop in which each upgraded system further improves the process that produces its successors.
OpenMLE studies a concrete step from evolution toward meta-evolution: executable search experience is reused through supervised fine-tuning and reinforcement learning to improve the program-transformation model.

\textbf{Executable evolution for AI-building tasks.}
We use machine learning engineering (MLE) as a measurable testbed for this process.
Following MLE-Dojo and MLE-Bench~\citep{qiang2025mledojo,chan2024mlebench}, each task $\tau$ contains a natural-language specification, visible data assets, a submission contract, a task-specific evaluator, and a sandboxed execution environment.
At step $t$, the inference-time search algorithm selects an operator
$a_t$ and constructs its operator context $c_t$ from zero or more parent
programs and their execution feedback.
The model then proposes
\[
p_t \sim g_\theta(\cdot\mid\tau,a_t,c_t),
\qquad
s_t=R_\tau(\mathcal{E}(p_t,\tau)).
\]
Here, $g_\theta$ denotes the operator-conditioned program-generation
policy parameterized by the language-model parameters $\theta$.
Given task $\tau$, operator $a_t$, and context $c_t$, it defines a
distribution over candidate programs.
The sandbox $\mathcal{E}$ executes $p_t$ and returns the task score $s_t$ mapped by a task-specific evaluator $R_\tau$ together
with status, logs, artifacts, and runtime metadata.
Because task metrics have different ranges and directions, Section~\ref{sec:execution-grounded-rl} converts $s_t$ into a signed score $\tilde{s}_t$, for which larger is always better, and then into a normalized reward.

Each program and its execution feedback are stored in a task-local
program database, and program-transformation operators are repeatedly
applied to generate new candidates.
Within a finite execution budget, evolutionary inference seeks the
candidate with the highest signed score:
\[
p^\star
=
p_{\arg\max_{t\in\mathcal{I}}\tilde{s}_t},
\]
where $\mathcal{I}$ denotes the index set of all candidate programs
recorded in the task-local program database during the budgeted search.
OpenMLE instantiates the operator space with \textsc{Draft},
\textsc{Improve}, \textsc{Debug}, and \textsc{Crossover}, while leaving
their global composition to the inference-time search algorithm.

\textbf{Learning to evolve from executable experience.}
Meta-evolution improves $\theta$ so that
$g_\theta(\cdot\mid\tau,a,c)$ assigns higher probability to programs
with stronger execution outcomes.
Both SFT and RL can be summarized as optimizing
\[
\mathcal{L}_{\mathrm{evo}}(\theta)
=
-
\mathbb{E}_{(\tau_i,a_i,c_i,p_i)}
\left[
w(s_i)
\log g_\theta
\left(
p_i\mid\tau_i,a_i,c_i
\right)
\right],
\]
where
\[
s_i=R_{\tau_i}\!\left(\mathcal{E}(p_i,\tau_i)\right)
\]
is the execution score and $w(s_i)$ converts this outcome into a
learning weight.
For SFT, quality filtering retains high-scoring programs $p_i^+$ and
assigns them positive supervision.
For RL, newly sampled programs are weighted by their processed
execution rewards and entropic advantages within the clipped policy
objective.
Thus, both stages update the same parameters $\theta$ to make
high-quality executable programs more likely under
$g_\theta(\cdot\mid\tau,a,c)$.
